% Figura: perceptrón multicapa, pasada hacia adelante y retropropagación.
% Se usa en sec:marco_fundamentos, subsección de entrenamiento.

\begin{figure}[htbp]
    \centering
    \begin{tikzpicture}[x=1mm, y=1mm,
        neurona/.style = {circle, draw=figtrazo, minimum size=6mm, inner sep=0pt, font=\footnotesize},
    ]

        % Capas.
        \foreach \i in {1, 2, 3} {
            \node[neurona, fill=figazul!15, draw=figazul!80!black] (x\i) at (0, 16 - \i*8) {$x_{\i}$};
        }
        \foreach \i in {1, 2, 3, 4} {
            \node[neurona, fill=fignaranja!18, draw=fignaranja!70!black] (h\i) at (25, 20 - \i*8) {};
        }
        \foreach \i in {1, 2, 3} {
            \node[neurona, fill=figverde!12, draw=figverde!70!black] (z\i) at (50, 16 - \i*8) {$z_{\i}$};
        }

        % Conexiones.
        \foreach \i in {1, 2, 3} { \foreach \j in {1, 2, 3, 4} {
            \draw[figtrazo!60, thin] (x\i) -- (h\j);
        } }
        \foreach \j in {1, 2, 3, 4} { \foreach \k in {1, 2, 3} {
            \draw[figtrazo!60, thin] (h\j) -- (z\k);
        } }

        % Softmax y pérdida.
        \node[figtarea, text width=16mm, minimum height=8mm] (sm) at (72, 0) {softmax};
        \node[figsalida, text width=20mm, minimum height=8mm] (loss) at (100, 0)
            {pérdida\\$\mathcal{L} = -\log p_y$};
        \draw[figflecha] (56, 0) -- (sm.west);
        \draw[figflecha] (sm.east) -- node[figetiqueta, above] {$p_c$} (loss.west);

        % Rótulos de capas.
        \node[font=\footnotesize, text=figtrazo] at (0, -22) {entrada};
        \node[font=\footnotesize, text=figtrazo, align=center] at (25, -22) {capa oculta\\$\phi(W\mathbf{x}+\mathbf{b})$};
        \node[font=\footnotesize, text=figtrazo] at (50, -22) {\emph{logits}};

        % Pasadas.
        \draw[figflecha, figazul!80!black, thick] (0, 24) -- node[figetiqueta, above, text=figazul!80!black]
            {pasada hacia adelante: se calcula la salida} (100, 24);
        \draw[-{Latex[length=2mm, width=1.6mm]}, thick, densely dashed, figbermellon]
            (100, -32) -- node[figetiqueta, below, text=figbermellon]
            {retropropagación: se calcula $\nabla_\theta \mathcal{L}$ capa por capa} (0, -32);

    \end{tikzpicture}
    \caption[Perceptrón multicapa y retropropagación]{Perceptrón multicapa con una capa oculta. La pasada
    hacia adelante transforma la entrada capa por capa hasta los \emph{logits}, que la softmax convierte en
    probabilidades, y la pérdida mide la probabilidad asignada a la clase correcta. La retropropagación
    recorre el camino inverso con la regla de la cadena y obtiene el gradiente de la pérdida respecto de
    cada parámetro.}
    \label{fig:marco_perceptron}
\end{figure}
